\section{Introduction}
\label{sec:introduction}

Portugal introduced a statutory minimum wage in May 1974, in the weeks after
the revolution, at 3\,300 escudos a month. Over the
\ObservedYears{}~years to \MacroLastYear{} the nominal figure rose by a factor
of more than fifty. Its real value tells a different story, and that story is
the subject of this paper.

We assemble the statutory history from primary law -- every act retrieved from
the Diário da República and read for the value it sets -- and combine it with a
consumer price index and a productivity series that both reach back to the
1960s. Three findings follow.

The first is a matter of arithmetic, and it is the paper's central fact. Indexed
to \MacroFirstYear{}, labour productivity reached \ProductivityIndexEnd{} by
\MacroLastYear{} while the real minimum wage reached \RealWageIndexEnd{}. The
real wage floor fell through the inflation of the late 1970s and early 1980s and
did not regain its introductory level for roughly four decades. Measured against
output per worker, it bottomed at \WageToProductivityMin{} in
\WageToProductivityMinYear{} and stood at \WageToProductivityEnd{} at the end of
the sample: an hour of minimum-wage work commands about half the share of
national output it did when the floor was created.

The second concerns policy rather than outcomes. Decomposing statutory growth
against a benchmark that compounds productivity growth with the previous year's
inflation yields a residual that averages \MeanResidualPct{}~percentage points a
year and cumulates to \CumulativeResidualPct{}~points. That average conceals a
reversal: the residual is strongly negative through the revolutionary
transition, close to zero from the late 1990s, and turns persistently positive
from 2015. The Portuguese minimum wage spent most of its life losing ground to
the benchmark and has spent the last decade regaining it.

The third finding is negative, and we state it as a result rather than as a
limitation. The natural way to identify the effect of minimum-wage policy on
prices in Portugal is to exploit the autonomous regions, which may set a wage
above the mainland's. We show that this cannot work with the variation that
exists. The Azores legislate a permanent proportional supplement, so after the
month it was introduced their statutory change is identical to the mainland's in
logs and contributes no independent timing at all. Madeira legislates a value,
and is the only source of genuine divergence. With the small number of
Portuguese regions, cluster-robust inference rejects a true null far above its
nominal size, and the estimates that conventional standard errors call highly
significant survive none of a wild cluster bootstrap.

We therefore make no causal claim about pass-through, and explain in
Section~\ref{sec:robustness} what would be required to support one. What the
paper offers instead is a measured long-run account of a policy whose real
generosity has moved a great deal, and the reproducible statutory series needed
to study it.

\paragraph{Contribution.} The statutory series itself is a contribution. Prior
work on Portuguese minimum wages relies on secondary compilations that flatten
features of the law: that three legally distinct minimum wages coexisted until
1991 and two until 2004, that the published official history omits the act
effective in 2000, and that the two autonomous regions set their wages by
different mechanisms. Each is documented here and each is read from the act
itself. The series, the regional schedules and the pipeline that builds them are
released with the paper.
